\section{The MLIPs Knowledge Graph}
\label{sec:kg}
%% ====================================================================

The \MLIPsKG{} is a structured knowledge graph grounded in the \MLIPsOntology{}
and populated in collaboration with domain experts from the META-LEARN project.

\subsection{Scope and Content}

The knowledge graph is organized around six entity types:

\begin{enumerate}
\item \emph{MLIP algorithms:} MACE, ACE, HDNNP, NequIP, M3GNet, MTP, and
  others, with their architectural characteristics and hyperparameter spaces.

\item \emph{Libraries and tools:} ASE~\cite{larsen2017ase}, LAMMPS, VASP,
  GPAW, and MLIP-specific training frameworks, with version information and
  interoperability notes.

\item \emph{Materials and systems:} Elements, alloys, and compound classes,
  including microstructural complexity (point defects, dislocations, surfaces,
  grain boundaries).

\item \emph{Simulation types:} Molecular dynamics, Monte Carlo, geometry
  optimization, phonon calculations, and thermodynamic integration.

\item \emph{Training datasets:} Provenance, size, property coverage, and DFT
  settings used to generate training data.

\item \emph{Papers and benchmarks:} Published studies linked to the
  algorithms, materials, and accuracy metrics they report.
\end{enumerate}

Where possible, entities are linked to their corresponding
Wikidata~\cite{vrandecic2014wikidata} items, connecting the \MLIPsKG{} to the
broader Linked Data ecosystem and enabling federated queries.

%% TODO: add statistics table (number of instances per class, total triples)

\subsection{Construction Methodology}

The knowledge graph is constructed through three complementary channels:

\begin{enumerate}
\item \emph{Expert curation.} Domain experts from the META-LEARN project's
  seven work packages contribute structured data through the SHACL-driven
  web editor, ensuring accuracy of domain-specific details (hyperparameter
  values, DFT settings, accuracy metrics).

\item \emph{AI-assisted extraction.} A specialized AI agent reads papers from
  the team's literature database and extracts structured statements about
  algorithms, training datasets, hyperparameters, and reported results. These
  statements are stored in named graphs with provenance metadata (source paper,
  extraction method, confidence level) and are curated by domain experts before
  incorporation.

\item \emph{Programmatic population via \SemASE{}.} Simulation data from
  ASE-based workflows is serialized to RDF and ingested via SPARQL UPDATE.
\end{enumerate}

\subsection{Maintenance and Versioning}

The knowledge graph is publicly queryable via a SPARQL endpoint. Periodic RDF
dumps are archived on DaRUS~\cite{darus2024} (the University of Stuttgart's
Dataverse repository) for long-term preservation with persistent DOIs. The
ontology is versioned alongside the knowledge graph, with change logs
documenting additions and modifications.

%% ====================================================================
